\begin{idstatement}{C.2.INT.3}{Proved sign and equality for every reducible fibre}
\label{st:C.2.INT.3}
Let $f:Y\to B$ satisfy the hypotheses of C.2.INT.2,
and let $S$ be a connected reducible fibre support.
Then
\[
 e(S)\ge0.
 \tag{C.INT.15}
\]
Equality holds precisely when the normalization of $S$
is a disjoint union of smooth minimal Hopf or
non-K\"ahler elliptic surfaces, and the reduced
conductor image and its reduced inverse image are
disjoint unions of smooth elliptic curves. The maps
between these elliptic curves are finite unramified
covers. This criterion allows arbitrary conductor
scheme structure and arbitrary unequal multiplicities
of the full fibre.
\end{idstatement}
\begin{proof}
Let $\nu:T=\coprod T_i\to S$ be normalization,
$\pi_i:Z_i\to T_i$ its minimal resolutions, and
$\sigma_i:Z_i\to Z_i^0$ the minimal-model maps.
By C.2.INT.1--C.2.INT.2 the $Z_i^0$ are Hopf or
non-K\"ahler elliptic surfaces. They have no rational
curves and have both Euler characteristics
$e(Z_i^0)=\chi(\OO_{Z_i^0})=0$.
Every rational curve of $Z_i$ is exceptional for
$\sigma_i$; their total number is exactly the number
$N_i$ of point blowups, and $e(Z_i)=N_i$.
The elliptic strict transforms cannot be contracted
by $\pi_i$, by C.2.INT.1. Thus every
$\pi_i$-exceptional curve is $\sigma_i$-exceptional.

The connected fibres of $\pi_i$ consequently map to
points under $\sigma_i$. Normality and
$(\pi_i)_*\OO_{Z_i}=\OO_{T_i}$ descend the latter
map to a proper birational morphism
$\tau_i:T_i\to Z_i^0$, with
\[
 \sigma_i=\tau_i\pi_i.
 \tag{C.INT.16}
\]
This is an actual factorization, not a replacement
of the normalization by its minimal model.
Since point blowups have $R\sigma_{i*}\OO_{Z_i}
=\OO_{Z_i^0}$, the low-degree Leray sequence gives
\[
 0\longrightarrow R^1\tau_{i*}\OO_{T_i}
 \longrightarrow 0
 \longrightarrow \tau_{i*}R^1\pi_{i*}\OO_{Z_i}
 \longrightarrow R^2\tau_{i*}\OO_{T_i}=0.
 \tag{C.INT.17}
\]
The last vanishing follows from fibre dimension at
most one. The middle sheaf $R^1\pi_{i*}\OO_{Z_i}$
is supported at finitely many points, and its
pushforward under $\tau_i$ detects it. Hence it is
zero. This proves that $T_i$ has rational
singularities, and in particular
$\chi(\OO_{T_i})=\chi(\OO_{Z_i})=0$.

Let $\Gamma$ be the full conductor image scheme,
$A_i\subset T_i$ its full inverse-image conductor,
and let $D_i\ge0$ be the resolution conductor from
A.5.3. Its line relation
$K_{Z_i}+D_i=h_i^*(K_Y+S)$ has torsion topological
first Chern class. That line descends through
$\sigma_i$. Consequently the exact canonical
divisor formula gives
\[
 D_i=\sigma_i^*D_i^0-K_{Z_i/Z_i^0},
 \qquad D_i^0=(\sigma_i)_*D_i.
 \tag{C.INT.18}
\]
The divisor $D_i^0$ is nonzero, since $D_i$ contains
a nonzero adjacent component divisor whose
minimal-model pushforward is nonzero by (C.INT.2).
It is a sum of positive multiples of disjoint
smooth elliptic curves. At each point blowup,
(C.INT.18) has coefficient
$\operatorname{mult}_pD_{\rm old}-1$ on the new
exceptional curve. Effectivity makes the centre
lie on the old support. Starting from the disjoint
elliptic curves, the reduced support consequently
acquires only rational leaves and subdivisions of
edges. Each of its connected components is one
elliptic curve with rational trees attached.
Writing $r_i$ for its number of rational curves,
we obtain
\[
 e((D_i)_{\rm red})=r_i.
 \tag{C.INT.19}
\]

We next prove a simultaneous cohomological
comparison on the entire, possibly nonreduced,
conductor. Relative duality applied to
$R\pi_{i*}\OO_{Z_i}=\OO_{T_i}$ gives
$R\pi_{i*}\omega_{Z_i}=\omega_{T_i}$.
Finite conductor duality identifies
$\omega_{T_i}\otimes\nu^*\omega_S^{-1}$ with
the conductor ideal $\mathcal I_{A_i/T_i}$.
The divisor sequence on $Z_i$, using
$\OO_{Z_i}(-D_i)=\omega_{Z_i}\otimes h_i^*\omega_S^{-1}$,
therefore gives
\[
 R\pi_{i*}\OO_{D_i}=\OO_{A_i}.
 \tag{C.INT.20}
\]
The identical calculation for $\sigma_i$, using
(C.INT.18), gives
\[
 R\sigma_{i*}\OO_{D_i}=\OO_{D_i^0},
 \qquad R\tau_{i*}\OO_{A_i}=\OO_{D_i^0}.
 \tag{C.INT.21}
\]
For completeness the canonical pushforward for a
point blowup is elementary: its relative canonical
line is $\OO(E)$, and
$0\to\OO\to\OO(E)\to\OO_E(-1)\to0$
has zero cohomology in its last term. It gives
$R\sigma_*\omega=\omega$; composition proves the
version used in (C.INT.21).

For $M\in\Pic^0(Y)$ its pullback to $Z_i$ descends
to a line $M_i^0\in\Pic^0(Z_i^0)$, by the Picard
blowup isomorphism. Projection formula and
(C.INT.20)--(C.INT.21) identify the actual spaces
\[
 H^0(A_i,\nu^*M|_{A_i})
   =H^0(D_i^0,M_i^0|_{D_i^0}).
 \tag{C.INT.22}
\]
For a component $nE$ of $D_i^0$, the filtration by
its powers of the defining ideal has graded lines
$M_i^0|_E\otimes\OO_E(-jE)$, $0\le j<n$.
They have degree zero, since $E^2=0$ on the
minimal model. C.2.INT.1 proves that the map from
$\Pic^0(Y)$ to $\Pic^0(E)$ is nonconstant. We can
therefore choose one $M$ avoiding, for each of
the finitely many such curves and integers $j$,
the prescribed value $\OO_E(jE)$.
Each excluded locus is a proper closed analytic
subset of the connected complex Lie group
$\Pic^0(Y)$; their finite union does not fill it.
All the graded lines are nontrivial of degree
zero on an elliptic curve and have no sections.
The filtration and (C.INT.22) prove
\[
 H^0(A_i,\nu^*M)=0\quad\hbox{for every }i.
 \tag{C.INT.23}
\]

The conductor inclusion is
$\OO_\Gamma\hookrightarrow\nu_*\OO_A$, where
$A=\coprod A_i$. It now gives $H^0(\Gamma,M)=0$.
The normalization sequence and the preceding
rational-singularity computation give
$\chi(\mathcal Q)=0$ for
$\mathcal Q=\nu_*\OO_T/\OO_S$: here
$\chi(\OO_S)=0$ is the Cartier-divisor
Riemann--Roch calculation under $H^2(Y,\Q)=0$.
Conductor duality is
$\mathcal Q=\omega_\Gamma\otimes\omega_S^{-1}|_\Gamma$.
Twisting on this compact Cohen--Macaulay curve by
an ambient topologically torsion line preserves
Euler characteristic; its degree on every
component is zero. Serre duality thus gives
\[
 \chi(\OO_\Gamma)=-\chi(\mathcal Q)=0.
 \tag{C.INT.24}
\]
The same holds after twisting by $M$, so
$H^0(\Gamma,M)=0$ implies $H^1(\Gamma,M)=0$.
Restriction to $C=\Gamma_{\rm red}$ surjects
on $H^1$, since its kernel is supported in
dimension at most one and has no $H^2$.
Consequently
\[
 \chi(\OO_C)=\chi(M|_C)=h^0(C,M|_C)\ge0.
 \tag{C.INT.25}
\]
This is where the full conductor scheme is used;
no passage to its reduction is assumed harmless.

Each normalized component of $C$ has genus zero
or one, since it is dominated by a curve of
$(D_i)_{\rm red}$. Let $n_0$ be the number of
rational-normalized components, and let
$\delta_x$ and $s_x$ be the delta invariant and
branch number of a singular point of $C$.
The exact curve normalization formulas give
\[
 e(C)=\chi(\OO_C)+n_0+
                \sum_x\bigl(\delta_x-(s_x-1)\bigr)\ge0.
 \tag{C.INT.26}
\]
Every summand is nonnegative, because
$\delta_x\ge s_x-1$. Equality holds precisely
when $C$ is a disjoint union of smooth elliptic
curves. Indeed equality first forces $n_0=0$;
then $\chi(\OO_C)=-\sum\delta_x\ge0$
forces every $\delta_x=0$. The converse is immediate.

We finish the topological comparison on each
normalization component. A $\pi_i$-exceptional
curve outside $(D_i)_{\rm red}$ is a $(-2)$ curve
disjoint from that support, by the adjunction
argument following (C.INT.10). Each connected
configuration of such curves is a tree: a cycle
or a repeated intersection would contradict
negative definiteness of its matrix with
all diagonal entries $-2$. Such fibres lie
off $A_i$, because a fibre over its conductor
meets the strict conductor and is connected.
Let $a_i$ count their irreducible curves, and
let $b_i$ count exceptional fibres contained
in $(D_i)_{\rm red}$ whose image lies outside
$A_i$. Removing the conductor and contracting
these specified fibres gives the exact equality
\[
 e(T_i\setminus(A_i)_{\rm red})
       =N_i-r_i-a_i+b_i\ge0.
 \tag{C.INT.27}
\]
A tree of $a$ rational curves has Euler number
$a+1$, so its contraction removes $a$; a
positive-conductor fibre removed upstairs but
represented by a point downstairs adds one.
These are all changes, since normalization is
an isomorphism off the conductor. Nonnegativity
follows because the $r_i$ support curves and
the $a_i$ disjoint exceptional curves are
distinct rational curves among the $N_i$
rational curves of $Z_i$.
The normalization is an isomorphism
$T\setminus A\simeq S\setminus C$. Additivity
with (C.INT.26)--(C.INT.27) proves (C.INT.15).

If $e(S)=0$, every nonnegative summand in these
formulas is zero. In particular $C$ is a
union of smooth elliptic curves, $b_i=0$,
and $r_i+a_i=N_i$. Every rational curve in
$(D_i)_{\rm red}$ must then be $\pi_i$-exceptional:
a rational curve cannot map nonconstantly to
one of these elliptic curves. The $a_i$ curves
are exceptional by definition. Thus all $N_i$
rational curves of $Z_i$ are $\pi_i$-exceptional.
If $N_i>0$, the last point blowup in
$\sigma_i$ supplies a $(-1)$ curve on $Z_i$.
Its being $\pi_i$-exceptional contradicts the
minimality of that resolution. Hence $N_i=0$.
There are now no rational curves to contract,
and C.2.INT.1 excludes contracting elliptic
curves. Therefore $Z_i=T_i=Z_i^0$ is smooth
and minimal. Its reduced conductor is a
disjoint union of elliptic curves. A finite
map between smooth elliptic curves is
unramified by Riemann--Hurwitz, proving the
stated map comparison.
Conversely these normalization and conductor
conditions give $e(T_i)=e((A_i)_{\rm red})=e(C)=0$;
the normalization formula gives $e(S)=0$.
The underlying spaces, and all comparisons
just used, are independent of the positive
full-fibre multiplicities. This completes the
proof including the equality criterion.
\end{proof}


% C.2.IRR.1 from part_C2_irrational.tex

